\section{Talking Accordion Blues}

As a young boy of around 6 or 7, my grandparents expected me to play the accordion. Of course, they found a teacher who showed me all the important songs: ``Return to Sorrento'', ``Santa Lucia'', ``Malafemmina'', among others. They way it works is that you start with a small instrument --- I think mine had 24 buttons --- and then move on up to larger and larger accordions.

After less than a year, I was getting proficient and asked when I could get a new accordion. That's when I got the bad news that they discovered a congenital hernia that required surgery. My father explained that I wouldn't be able to lift a heavier instrument for a while, at least until the surgery healed, so my mother suggested I take piano lessons in the meantime. So I learned to play piano on an always out of tune upright, while waiting for my hernia to heal.

Eventually we moved to a larger house, so naturally we got a larger piano --- a sweet baby grand. But, no, there was no larger accordion waiting for me. As time passed, my interests changed and I forgot about the accordion.

This is a picture I took of a wandering accordion player in a cafe in Montmartre. A nice gig if you can get it --- fresh air, exercise, art, freedom ... he looks healthy enough, but I forgot to ask him if he had ever had hernia surgery. I know a few of you may be snickering, but he gets paid in euros and I rely on dollars.

It's been said that ``A gentleman is a man who can play the accordion ... and doesn't.''

I'm no gentleman.

Connie Francis got her start on the accordion, but blocked all youtube evidence. As she claimed:

\begin{quotex}
I was never encouraged to do it and I played the accordion, which I hated. I wish I had taken piano because I definitely would have written more songs of my own, but I didn't.
\end{quotex} 

\url{https://www.youtube.com/watch?v=0\_URJbCLqXE}

\flrightit{Posted on 2008-06-30 by Cologero}

